% ----------------------
\subsection{Section 128 (7 September 1842)}
% ----------------------
	\label{DC128}
	
	% -----
	\subsubsection{Historical Background}
	% -----
		\label{DC128-HB}
		
		% --
		\paragraph{JSP}
		% --
		
			``On 7 September 1842, JS dictated a letter addressed to the church giving further instructions on performing and recording baptisms for the dead. He wrote from the home of Edward Hunter, where he was hiding from men sent to arrest him on charges of complicity in the attempted murder of former Missouri governor Lilburn W. Boggs. This letter expanded on JS's 1 September 1842 letter to the church, which included some instructions on recording baptisms for the dead and promised to give additional instructions in the future. 

			``Since the founding of the church, record keeping had served as an important theme in both the revelations and teachings of JS and had influenced the day-to-day operations of the church. With the instructions in both this and the 1 September letter, however, JS outlined a more detail-oriented approach to record keeping than had previously been practiced among the Latter-day Saints. JS's discussion of record keeping in this 7 September letter was a vital step toward the procedural systematization of baptisms for the dead and the attendant records of those ordinances, as record keeping had been limited prior to September 1842 and often varied from branch to branch. The enhanced record-keeping efforts that JS suggested resembled those that contemporaneous state and county recorders were making to maintain official copies of deeds and other records. JS explained, however, that the Saints needed to maintain an accurate record of these ordinances because the record would subsequently be written in heaven and become a book that the Saints would offer to God. 

			``In the letter, JS also tied baptisms for the dead to his Nauvoo-era teachings regarding priesthood and dispensations of the gospel. He explained that through the ordinance, the Saints could forge a generational chain between parents and children, just as there would be a `welding together' of the various keys and dispensations from Adam down to JS in the `dispensation of the fulness of times.' Building on this teaching, JS used the letter to recount briefly that he had been visited by ancient prophets who provided him with the necessary keys. 

			``At JS's request, this letter was `read to the saints at the Grove near the Temple' on 11 September 1842. William Clayton reported in JS's journal that the letter `made a deep and solemn impression on the minds of the saints,' who `manifested their intentions to obey the instructions to the letter.' Soon thereafter, general church recorder James Sloan began keeping a new record for baptisms for the dead. The first entry in the book was for proxy baptisms performed on the evening of 11 September 1842, with the entries for that date reflecting the new instructions contained in this letter. 

			The featured version of the letter is a loose copy in Clayton's handwriting. It is the earliest known extant manuscript copy and may be the original dictated letter. The letter was copied into JS's journal by Eliza R. Snow on or around 11 September. Differences between the Clayton version and the version copied in JS's journal are noted. Although both versions date the letter 6 September, JS's journal entry for 7 September notes that he `wrote---or rather dictated a long Epistle to the Saints which he ordered to be read next Sabbath.' The letter may have been misdated in the loose copy, with the error copied into the journal, or Clayton may have mistakenly attributed the letter to 6 September while making entries in JS's journal. The letter was subsequently published in the 1 October 1842 issue of the Times and Seasons. At some point between September 1842 and July 1843, James Sloan inscribed onto a single leaf excerpts from the letter pertaining to how the records were to be kept. Water damage and other markings on the page suggest he may have kept this document with the baptisms for the dead records that he made pursuant to the instructions in the letter. The letter was later included in the 1844 edition of the Doctrine and Covenants.''
			
		% --
		\paragraph{HDDC}
		% --
		
			See HDDC 1669--1676 for more background to the introduction of vicarious ordinances and to this section of the DC. I have some of the material discussed there in my JS sources; see \hyperref[EMS-1840-JS-3-5-OCT-1840]{this discourse}, including the accounts by Vilate Kimball and Phebe Woodruff.
					
	% -----
	\subsubsection{Versions}
	% -----
		\label{DC128-V}
		
		See the \href{https://drive.google.com/file/d/1goAvNz8jS4R8slYfMG_db55Ty2z_vmgK/view?usp=sharing}{sections comparison} document.
	
		\begin{compactdesc}
			\item[JSP] \hyperref[EMS-1842-JS-7-SEP-1842]{7 September 1842}
			\item[BC] ---
			\item[DC 1835] ---
			\item[TS] 3:23 (1 October 1842), 934--936
			\item[DC 1844] Section 106, 420--430
			\item[DN] 5:43 (2 January 1856), 337
			\item[MS] 20:1 (2 January 1858), 5--8
		\end{compactdesc}

	% -----
	\subsubsection{Section 128:1} % *DC128-1 
	% -----
		\label{DC128-1}

		AS I stated to you in my letter before I left my place, that I would write to you from time to time and give you information in relation to many subjects, I now resume the subject of the baptism for the dead, as that subject seems to occupy my mind, and press itself upon my feelings the strongest, since I have been pursued by my enemies.

		% \begin{framed}
		% \scriptsize
			% \begin{compactdesc}
				% \item[JSP] 
				% % \item[BC] 
				% % \item[]
			% \end{compactdesc}
		% \end{framed}

	% -----
	\subsubsection{Section 128:2} % *DC128-2 
	% -----
		\label{DC128-2}

		I wrote a few words of revelation to you concerning a recorder. I have had a few additional views in relation to this matter, which I now certify. That is, it was declared in my former letter that there should be a recorder, who should be eye-witness, and also to hear with his ears, that he might make a record of a truth before the Lord.

		% \begin{framed}
		% \scriptsize
			% \begin{compactdesc}
				% \item[JSP] 
				% % \item[BC] 
				% % \item[]
			% \end{compactdesc}
		% \end{framed}

	% -----
	\subsubsection{Section 128:3} % *DC128-3 
	% -----
		\label{DC128-3}

		Now, in relation to this matter, it would be very difficult for one recorder to be present at all times, and to do all the business. To obviate this difficulty, there can be a recorder appointed in each ward of the city, who is well qualified for taking accurate minutes; and let him be very particular and precise in taking the whole proceedings,
			\footnote{%
				% \label{}%
				Note the nature of a ``witness' here for JS: someone who sees and hears something.
				}%
		certifying in his record that he saw with his eyes, and heard with his ears, giving the date, and names, and so forth, and the history of the whole transaction; naming also some three individuals that are present, if there be any present, who can at any time when called upon certify to the same, that in the mouth of two or three witnesses every word may be established.

		% \begin{framed}
		% \scriptsize
			% \begin{compactdesc}
				% \item[JSP] 
				% % \item[BC] 
				% % \item[]
			% \end{compactdesc}
		% \end{framed}

	% -----
	\subsubsection{Section 128:4} % *DC128-4 
	% -----
		\label{DC128-4}

		Then, let there be a general recorder, to whom these other records can be handed, being attended with certificates over their own signatures, certifying that the record they have made is true. Then the general church recorder can enter the record on the general church book, with the certificates and all the attending witnesses, with his own statement that he verily believes the above statement and records to be true, from his knowledge of the general character and appointment of those men by the church. And when this is done on the general church book, the record shall be just as holy, and shall answer the ordinance just the same as if he had seen with his eyes and heard with his ears, and made a record of the same on the general church book.

		% \begin{framed}
		% \scriptsize
			% \begin{compactdesc}
				% \item[JSP] 
				% % \item[BC] 
				% % \item[]
			% \end{compactdesc}
		% \end{framed}

	% -----
	\subsubsection{Section 128:5} % *DC128-5 
	% -----
		\label{DC128-5}

		You may think this order of things to be very particular; but let me tell you that it is only to answer the will of God, by conforming to the ordinance and preparation that the Lord ordained and prepared before the foundation of the world, for the salvation of the dead who should die without a knowledge of the gospel.

		% \begin{framed}
		% \scriptsize
			% \begin{compactdesc}
				% \item[JSP] 
				% % \item[BC] 
				% % \item[]
			% \end{compactdesc}
		% \end{framed}

	% -----
	\subsubsection{Section 128:6} % *DC128-6 
	% -----
		\label{DC128-6}

		And further, I want you to remember that John the Revelator was contemplating this very subject in relation to the dead, when he declared, as you will find recorded in \hyperref[REVELATIONNT-20-12]{Revelation 20:12}---\textit{And I saw the dead, small and great, stand before God; and the books were opened; and another book was opened, which is the book of life; and the dead were judged out of those things which were written in the books, according to their works.}

		% \begin{framed}
		% \scriptsize
			% \begin{compactdesc}
				% \item[JSP] 
				% % \item[BC] 
				% % \item[]
			% \end{compactdesc}
		% \end{framed}

	% -----
	\subsubsection{Section 128:7} % *DC128-7 
	% -----
		\label{DC128-7}

		You will discover in this quotation that the books were opened; and another book was opened, which was the book of life; but the dead were judged out of those things which were written in the books, according to their works; consequently, the books spoken of must be the books which contained the record of their works, and refer to the records which are kept on the earth. And the book which was the book of life is the record which is kept in heaven; the principle agreeing precisely with the doctrine which is commanded you in the revelation contained in the letter which I wrote to you previous to my leaving my place---that in all your recordings it may be recorded in heaven.

		% \begin{framed}
		% \scriptsize
			% \begin{compactdesc}
				% \item[JSP] 
				% % \item[BC] 
				% % \item[]
			% \end{compactdesc}
		% \end{framed}

	% -----
	\subsubsection{Section 128:8} % *DC128-8 
	% -----
		\label{DC128-8}

		Now, the nature of this ordinance consists in the power of the priesthood, by the revelation of Jesus Christ, wherein it is granted that whatsoever you bind on earth shall be bound in heaven, and whatsoever you loose on earth shall be loosed in heaven. Or, in other words, taking a different view of the translation, whatsoever you record on earth shall be recorded in heaven, and whatsoever you do not record on earth shall not be recorded in heaven; for out of the books shall your dead be judged, according to their own works, whether they themselves have attended to the ordinances in their own \textit{propria persona}, or by the means of their own agents, according to the ordinance which God has prepared for their salvation from before the foundation of the world, according to the records which they have kept concerning their dead.

		% \begin{framed}
		% \scriptsize
			% \begin{compactdesc}
				% \item[JSP] 
				% % \item[BC] 
				% % \item[]
			% \end{compactdesc}
		% \end{framed}

	% -----
	\subsubsection{Section 128:9} % *DC128-9 
	% -----
		\label{DC128-9}

		It may seem to some to be a very bold doctrine that we talk of---a power which records or binds on earth and binds in heaven. Nevertheless, in all ages of the world, whenever the Lord has given a dispensation of the priesthood to any man by actual revelation, or any set of men, this power has always been given. Hence, whatsoever those men did in authority, in the name of the Lord, and did it truly and faithfully, and kept a proper and faithful record of the same, it became a law on earth and in heaven, and could not be annulled, according to the decrees of the great Jehovah. This is a faithful saying. Who can hear it?

		% \begin{framed}
		% \scriptsize
			% \begin{compactdesc}
				% \item[JSP] 
				% % \item[BC] 
				% % \item[]
			% \end{compactdesc}
		% \end{framed}

	% -----
	\subsubsection{Section 128:10} % *DC128-10 
	% -----
		\label{DC128-10}

		And again, for the precedent, Matthew 16:18, 19: \textit{And I say also unto thee, That thou art Peter, and upon this rock I will build my church; and the gates of hell shall not prevail against it. And I will give unto thee the keys of the kingdom of heaven: and whatsoever thou shalt bind on earth shall be bound in heaven; and whatsoever thou shalt loose on earth shall be loosed in heaven.}

		% \begin{framed}
		% \scriptsize
			% \begin{compactdesc}
				% \item[JSP] 
				% % \item[BC] 
				% % \item[]
			% \end{compactdesc}
		% \end{framed}

	% -----
	\subsubsection{Section 128:11} % *DC128-11 
	% -----
		\label{DC128-11}

		Now the great and grand secret of the whole matter, and the \textit{summum bonum} of the whole subject that is lying before us, consists in obtaining the powers of the Holy Priesthood. For him to whom these keys are given there is no difficulty in obtaining a knowledge of facts in relation to the salvation of the children of men, both as well for the dead as for the living.

		% \begin{framed}
		% \scriptsize
			% \begin{compactdesc}
				% \item[JSP] 
				% % \item[BC] 
				% % \item[]
			% \end{compactdesc}
		% \end{framed}

	% -----
	\subsubsection{Section 128:12} % *DC128-12 
	% -----
		\label{DC128-12}

		Herein is glory and honor, and immortality and eternal life---The ordinance of baptism by water, to be immersed therein in order to answer to the likeness of the dead, that one principle might accord with the other; to be immersed in the water and come forth out of the water is in the likeness of the resurrection of the dead in coming forth out of their graves; hence, this ordinance was instituted to form a relationship with the ordinance of baptism for the dead, being in likeness of the dead.

		% \begin{framed}
		% \scriptsize
			% \begin{compactdesc}
				% \item[JSP] 
				% % \item[BC] 
				% % \item[]
			% \end{compactdesc}
		% \end{framed}

	% -----
	\subsubsection{Section 128:13} % *DC128-13 
	% -----
		\label{DC128-13}

		Consequently, the baptismal font was instituted as a similitude of the grave, and was commanded to be in a place underneath where the living are wont to assemble, to show forth the living and the dead, and that all things may have their likeness, and that they may accord one with another---that which is earthly conforming to that which is heavenly, as Paul hath declared, 1 Corinthians 15:46, 47, and 48:

		% \begin{framed}
		% \scriptsize
			% \begin{compactdesc}
				% \item[JSP] 
				% % \item[BC] 
				% % \item[]
			% \end{compactdesc}
		% \end{framed}

	% -----
	\subsubsection{Section 128:14} % *DC128-14 
	% -----
		\label{DC128-14}

		\textit{Howbeit that was not first which is spiritual, but that which is natural; and afterward that which is spiritual. The first man is of the earth, earthy; the second man is the Lord from heaven. As is the earthy, such are they also that are earthy; and as is the heavenly, such are they also that are heavenly.} And as are the records on the earth in relation to your dead, which are truly made out, so also are the records in heaven. This, therefore, is the sealing and binding power, and, in one sense of the word, the keys of the kingdom, which consist in the key of knowledge. 	

		% \begin{framed}
		% \scriptsize
			% \begin{compactdesc}
				% \item[JSP] 
				% % \item[BC] 
				% % \item[]
			% \end{compactdesc}
		% \end{framed}

	% -----
	\subsubsection{Section 128:15} % *DC128-15 
	% -----
		\label{DC128-15}

		And now, my dearly beloved brethren and sisters, let me assure you that these are principles in relation to the dead and the living that cannot be lightly passed over, as pertaining to our salvation. For their salvation is necessary and essential to our salvation, as Paul says concerning the fathers---that they without us cannot be made perfect---neither can we without our dead be made perfect.

		% \begin{framed}
		% \scriptsize
			% \begin{compactdesc}
				% \item[JSP] 
				% % \item[BC] 
				% % \item[]
			% \end{compactdesc}
		% \end{framed}

	% -----
	\subsubsection{Section 128:16} % *DC128-16 
	% -----
		\label{DC128-16}

		And now, in relation to the baptism for the dead, I will give you another quotation of Paul, 1 Corinthians 15:29: \textit{Else what shall they do which are baptized for the dead, if the dead rise not at all? Why are they then baptized for the dead?}

		% \begin{framed}
		% \scriptsize
			% \begin{compactdesc}
				% \item[JSP] 
				% % \item[BC] 
				% % \item[]
			% \end{compactdesc}
		% \end{framed}

	% -----
	\subsubsection{Section 128:17} % *DC128-17 
	% -----
		\label{DC128-17}

		And again, in connection with this quotation I will give you a quotation from one of the prophets, who had his eye fixed on the restoration of the priesthood, the glories to be revealed in the last days, and in an especial manner this most glorious of all subjects belonging to the everlasting gospel, namely, the baptism for the dead; for Malachi says, last chapter, verses 5th and 6th: \textit{Behold, I will send you Elijah the prophet before the coming of the great and dreadful day of the Lord: And he shall turn the heart of the fathers to the children, and the heart of the children to their fathers, lest I come and smite the earth with a curse.}

		% \begin{framed}
		% \scriptsize
			% \begin{compactdesc}
				% \item[JSP] 
				% % \item[BC] 
				% % \item[]
			% \end{compactdesc}
		% \end{framed}

	% -----
	\subsubsection{Section 128:18} % *DC128-18 
	% -----
		\label{DC128-18}

		I might have rendered a plainer translation to this, but it is sufficiently plain to suit my purpose as it stands. It is sufficient to know, in this case, that the earth will be smitten with a curse unless there is a welding link of some kind or other between the fathers and the children, upon some subject or other---and behold what is that subject? It is the baptism for the dead. For we without them cannot be made perfect; neither can they without us be made perfect. Neither can they nor we be made perfect without those who have died in the gospel also; for it is necessary in the ushering in of the dispensation of the fulness of times, which dispensation is now beginning to usher in, that a whole and complete and perfect union, and welding together of dispensations, and keys, and powers, and glories should take place, and be revealed from the days of Adam even to the present time. And not only this, but those things which never have been revealed from the foundation of the world, but have been kept hid from the wise and prudent, shall be revealed unto babes and sucklings in this, the dispensation of the fulness of times.

		% \begin{framed}
		% \scriptsize
			% \begin{compactdesc}
				% \item[JSP] 
				% % \item[BC] 
				% % \item[]
			% \end{compactdesc}
		% \end{framed}

	% -----
	\subsubsection{Section 128:19} % *DC128-19 
	% -----
		\label{DC128-19}

		Now, what do we hear in the gospel which we have received? A voice of gladness! A voice of mercy from heaven; and a voice of truth out of the earth; glad tidings for the dead; a voice of gladness for the living and the dead; glad tidings of great joy. How beautiful upon the mountains are the feet of those that bring glad tidings of good things, and that say unto Zion: Behold, thy God reigneth! As the dews%
			\footnote{%
				% \label{}%
				See \hyperref[DC121-45]{DC 121:45}.
				}
		of \hyperref[CARMEL]{Carmel}, so shall the knowledge of God descend upon them!

		% \begin{framed}
		% \scriptsize
			% \begin{compactdesc}
				% \item[JSP] 
				% % \item[BC] 
				% % \item[]
			% \end{compactdesc}
		% \end{framed}

	% -----
	\subsubsection{Section 128:20} % *DC128-20 
	% -----
		\label{DC128-20}

		And again, what do we hear? Glad tidings from Cumorah! Moroni, an angel from heaven, declaring the fulfilment of the prophets---the book to be revealed. A voice of the Lord in the wilderness of Fayette, Seneca county, declaring the three witnesses to bear record of the book! The voice of Michael on the banks of the Susquehanna, detecting the devil when he appeared as an angel of light!%
			\footnote{%
				% \label{}%
				Compare \hyperref[2CORINTHIANS-11-14]{2 Corinthians 11:14}.
				}
		The voice of Peter, James, and John in the wilderness between Harmony, Susquehanna county, and Colesville, Broome county, on the Susquehanna river, declaring themselves as possessing the keys of the kingdom, and of the dispensation of the fulness of times!

		% \begin{framed}
		% \scriptsize
			% \begin{compactdesc}
				% \item[JSP] 
				% % \item[BC] 
				% % \item[]
			% \end{compactdesc}
		% \end{framed}

	% -----
	\subsubsection{Section 128:21} % *DC128-21 
	% -----
		\label{DC128-21}

		And again, the voice of God in the chamber of old Father Whitmer, in Fayette, Seneca county, and at sundry times, and in divers places through all the travels and tribulations of this Church of Jesus Christ of Latter-day Saints! And the voice of Michael, the archangel; the voice of Gabriel, and of Raphael, and of divers angels, from Michael or Adam down to the present time, all declaring their dispensation, their rights, their keys, their honors, their majesty and glory, and the power of their priesthood; giving line upon line, precept upon precept; here a little, and there a little; giving us consolation by holding forth that which is to come, confirming our hope!

		% \begin{framed}
		% \scriptsize
			% \begin{compactdesc}
				% \item[JSP] 
				% % \item[BC] 
				% % \item[]
			% \end{compactdesc}
		% \end{framed}

	% -----
	\subsubsection{Section 128:22} % *DC128-22 
	% -----
		\label{DC128-22}

		Brethren, shall we not go on in so great a cause? Go forward and not backward. Courage, brethren; and on, on to the victory! Let your hearts rejoice, and be exceedingly glad. Let the earth break forth into singing. Let the dead speak forth anthems of eternal praise to the King Immanuel, who hath ordained, before the world was, that which would enable us to redeem them out of their prison; for the prisoners shall go free.

		% \begin{framed}
		% \scriptsize
			% \begin{compactdesc}
				% \item[JSP] 
				% % \item[BC] 
				% % \item[]
			% \end{compactdesc}
		% \end{framed}

	% -----
	\subsubsection{Section 128:23} % *DC128-23 
	% -----
		\label{DC128-23}

		Let the mountains shout for joy, and all ye valleys cry aloud; and all ye seas and dry lands tell the wonders of your Eternal King! And ye rivers, and brooks, and rills, flow down with gladness. Let the woods and all the trees of the field praise the Lord; and ye solid rocks weep for joy! And let the sun, moon, and the morning stars sing together, and let all the sons of God shout for joy! And let the eternal creations declare his name forever and ever! And again I say, how glorious is the voice we hear from heaven, proclaiming in our ears, glory, and salvation, and honor, and immortality, and eternal life; kingdoms, principalities, and powers!

		% \begin{framed}
		% \scriptsize
			% \begin{compactdesc}
				% \item[JSP] 
				% % \item[BC] 
				% % \item[]
			% \end{compactdesc}
		% \end{framed}

	% -----
	\subsubsection{Section 128:24} % *DC128-24 
	% -----
		\label{DC128-24}

		Behold, the great day of the Lord is at hand; and who can abide the day of his coming, and who can stand when he appeareth? For he is like a refiner's fire, and like fuller's soap; and he shall sit as a refiner and purifier of silver, and he shall purify the sons of Levi, and purge them as gold and silver, that they may offer unto the Lord an offering in righteousness. Let us, therefore, as a church and a people, and as Latter-day Saints, offer unto the Lord an offering in righteousness; and let us present in his holy temple, when it is finished, a book containing the records of our dead, which shall be worthy of all acceptation.

		% \begin{framed}
		% \scriptsize
			% \begin{compactdesc}
				% \item[JSP] 
				% % \item[BC] 
				% % \item[]
			% \end{compactdesc}
		% \end{framed}

	% -----
	\subsubsection{Section 128:25} % *DC128-25 
	% -----
		\label{DC128-25}

		Brethren, I have many things to say to you on the subject; but shall now close for the present, and continue the subject another time. I am, as ever, your humble servant and never deviating friend,
		
		JOSEPH SMITH.

		% \begin{framed}
		% \scriptsize
			% \begin{compactdesc}
				% \item[JSP] 
				% % \item[BC] 
				% % \item[]
			% \end{compactdesc}
		% \end{framed}